\chapter{模24计数器}

\section{应用统一框架}

\begin{enumerate}
  \item \textbf{模值N}：$N = 24$
  \item \textbf{寄存器位宽k}：$k = \lceil \log_2 24 \rceil = 5$（需要5位，最多表示31）
  \item \textbf{状态转移}：count $\leftarrow$ count + 1
  \item \textbf{回零条件}：count = 23(10111)时，强制归零
\end{enumerate}

\section{关键状态分析}

模24计数器的状态范围：00000(0) $\to$ 10111(23) $\to$ 00000(0)

\textbf{模24常与时钟应用相关}：24小时制时钟的小时位。

\begin{center}
\begin{tikzpicture}
  \node[draw, circle] (s0) at (0:2) {0};
  \node[draw, circle] (s6) at (90:2) {6};
  \node[draw, circle] (s12) at (180:2) {12};
  \node[draw, circle] (s18) at (270:2) {18};

  \draw[->] (s0) arc (0:85:2);
  \draw[->] (s6) arc (90:175:2);
  \draw[->] (s12) arc (180:265:2);
  \draw[->] (s18) arc (270:355:2);

  \node at (0,0) {24状态};
  \node at (0,-3) {\small 0→1→...→22→23→0};
\end{tikzpicture}
\end{center}

\section{为什么是5位}

$2^4 = 16 < 24 < 2^5 = 32$，所以需要5位寄存器。

5位二进制最大值11111 = 31，但计数器只用到0-23（24个状态）。状态24-31(11000-11111)永远不会出现。

\begin{examtip}[位宽确定公式]
对于模N计数器，所需寄存器位数 = $\lceil \log_2 N \rceil$

\begin{itemize}
  \item 模16 → 4位（恰好用完）
  \item 模12 → 4位（有4个冗余状态）
  \item 模24 → 5位（有8个冗余状态）
  \item 模60 → 6位（有4个冗余状态）
  \item 模10 → 4位（有6个冗余状态）
\end{itemize}
\end{examtip}

\section{VHDL实现}

\begin{lstlisting}[caption={模24计数器}]
library ieee;
use ieee.std_logic_1164.all;
use ieee.std_logic_unsigned.all;

entity counter_mod24 is
  port (
    clk   : in  std_logic;
    rst_n : in  std_logic;
    count : out std_logic_vector(4 downto 0);
    carry : out std_logic
  );
end counter_mod24;

architecture behavioral of counter_mod24 is
  signal cnt : std_logic_vector(4 downto 0);
begin
  process(clk, rst_n)
  begin
    if rst_n = '0' then
      cnt <= (others => '0');
    elsif rising_edge(clk) then
      if cnt = "10111" then    -- = 23
        cnt <= (others => '0');
      else
        cnt <= cnt + 1;
      end if;
    end if;
  end process;

  count <= cnt;
  carry <= '1' when cnt = "10111" else '0';
end behavioral;
\end{lstlisting}

\section{与模12的比较}

\begin{center}
\begin{tabular}{|l|c|c|}
\hline
\textbf{参数} & \textbf{模12} & \textbf{模24} \\
\hline
模值N & 12 & 24 \\
位宽 & 4位 & 5位 \\
清零条件 & cnt = 11(1011) & cnt = 23(10111) \\
其余逻辑 & \multicolumn{2}{c|}{\textbf{完全相同}} \\
\hline
\end{tabular}
\end{center}

\textbf{结论}：模12和模24的VHDL代码，区别\textbf{只有一行}——清零条件不同。这就是统一框架的力量。
